\chapter{Du fichier au graphe : le chargement}

Entre le moment où vous tapez \texttt{jenga build} et celui où le compilateur
démarre, Jenga fait un travail qu'on ne voit pas. Ce chapitre le décrit, parce
que la moitié des surprises viennent de là.

\section{Les six temps du chargement}

\begin{center}
\begin{tabular}{@{}rll@{}}
\toprule
& \textbf{Temps} & \textbf{Ce qui se passe} \\
\midrule
1 & trouver le fichier d'entrée & le \texttt{.jenga} de la racine \\
2 & préparer l'environnement & namespace, chemin des \emph{typings} \\
3 & se placer & le répertoire courant devient celui du fichier \\
4 & \textbf{exécuter} & \texttt{exec(\dots)} --- vos projets naissent ici \\
5 & récupérer le workspace & \texttt{getcurrentworkspace()} \\
6 & post-traiter & résolution des chemins, vérifications \\
\bottomrule
\end{tabular}
\end{center}

\begin{nkcode}{Jenga/Core/Loader.py, les lignes qui comptent}
old_cwd = Path.cwd()
os.chdir(filePath.parent)
...
_loadedFiles.append(filePath)
exec(filePath.read_text(encoding='utf-8-sig'), globals_dict)
workspace = Api.getcurrentworkspace()
if workspace is None:
    raise RuntimeError("No workspace defined in the entry file.")
self._PostProcessWorkspace(workspace, filePath)
\end{nkcode}

\begin{nkretenir}
\textbf{\texttt{utf-8-sig}, et non \texttt{utf-8}.} Le suffixe \texttt{-sig} fait
ignorer la marque d'ordre d'octets que certains éditeurs Windows placent en tête
de fichier.

Sans cela, un \texttt{.jenga} enregistré par le Bloc-notes échouerait sur sa
première ligne avec une erreur de syntaxe portant sur un caractère invisible.
C'est trois lettres, et elles évitent une demi-journée.
\end{nkretenir}

\section{Les projets naissent en s'exécutant}

\begin{nkretenir}
\textbf{Il n'y a pas d'étape « analyse ».} Quand \texttt{exec} se termine, le
graphe est déjà là : chaque \texttt{with project(\dots)} a créé un objet, chaque
\texttt{files(\dots)} a rempli une liste.

C'est pourquoi l'ordre d'écriture compte : une fonction doit être définie
\emph{avant} le \texttt{include} qui s'en sert, comme dans n'importe quel script.
\end{nkretenir}

\begin{nkcode}{Jenga/Core/Api.py -- ce que fait `with project`}
class project:
    def __init__(self, name: str):
        ...
\end{nkcode}

\begin{nkretenir}
\textbf{Les fonctions du vocabulaire écrivent dans une variable globale.}
\texttt{files()} ne reçoit pas de projet : il lit \texttt{\_currentProject}, que
\texttt{with project} vient de poser.

C'est simple, c'est efficace, et cela explique le piège du chapitre 4 : un appel
hors bloc écrit dans le \emph{dernier} projet, ou dans le vide.
\end{nkretenir}

\section{Le post-traitement}

\begin{nkretenir}
Après l'exécution, Jenga résout ce qui était relatif : les chemins deviennent
absolus, les motifs de fichiers sont développés, les dépendances sont mises en
correspondance avec des projets réels.

\textbf{C'est ici que la plupart des erreurs de description apparaissent} --- et
elles apparaissent \emph{loin} de la ligne fautive, puisque cette ligne s'est
exécutée sans se plaindre.
\end{nkretenir}

\section{Que fait Jenga quand ça rate}

\begin{nkcode}{Jenga/Core/Loader.py}
except Exception as e:
    Colored.PrintError(f"Error loading workspace: {e}")
    if self.verbose:
        traceback.print_exc()
    return None
finally:
    os.chdir(old_cwd)
\end{nkcode}

\begin{nkretenir}
\textbf{Trois choses à retenir de ces sept lignes.}

\textbf{Le \texttt{finally} restaure le répertoire courant} même en cas d'échec.
Sans lui, une erreur de chargement vous laisserait dans un autre dossier que
celui d'où vous avez tapé la commande.

\textbf{\texttt{-{}-verbose} donne la pile complète.} C'est le premier réflexe : le
message seul dit \emph{quoi}, la pile dit \emph{où}.

\textbf{Le chargeur rend \texttt{None}}, il ne lève pas. L'appelant doit donc
tester --- et c'est le genre de contrat qu'on oublie en écrivant un outil autour
de Jenga.
\end{nkretenir}

\section{Le namespace propagé}

\begin{nkcode}{Jenga/Core/Loader.py}
# Exposer le namespace LIVE du workspace racine pour que `include()` propage
# automatiquement aux fichiers inclus tout ce que l'utilisateur definit ici
# (config/fonctions/classes/imports) -> plus besoin de les re-importer dans
# chaque .jenga inclus.
Api._workspaceGlobals = globals_dict
\end{nkcode}

\begin{nkretenir}
\textbf{Ce qui est défini dans le fichier racine est visible partout.} C'est ce
qui permet à trente descripteurs d'application d'employer
\texttt{nkentseudependson} sans jamais l'importer.

Et le nettoyage est fait : \texttt{Api.\_workspaceGlobals = None} en fin de
chargement, « évite toute fuite de namespace entre deux chargements de
workspace ».
\end{nkretenir}

\begin{nkpiege}
\textbf{Cette commodité a un coût de lisibilité.} Un descripteur d'application
appelle une fonction dont rien, dans son fichier, ne dit d'où elle vient. Ni
\texttt{import}, ni définition.

Quand vous lisez un \texttt{.jenga} et rencontrez un nom inconnu, le réflexe est
donc : \textbf{chercher dans le fichier racine et dans ce qu'il inclut avant
lui}, jamais dans la documentation de Jenga.
\end{nkpiege}

\section{Le fichier d'entrée n'est pas toujours celui qu'on croit}

\begin{nkretenir}
Jenga cherche le \texttt{.jenga} en remontant depuis le répertoire courant. Vous
pouvez le forcer :

\begin{nkcode}{}
jenga info --jenga-file chemin/vers/Autre.jenga
\end{nkcode}

\textbf{Lancer une commande depuis un sous-dossier peut donc charger un workspace
différent} de celui que vous croyez. La sortie de \texttt{jenga info} donne
l'\texttt{Entry file} en toutes lettres --- c'est la ligne qui tranche.
\end{nkretenir}

\begin{nkbilan}
Le chargement tient en six temps, et le quatrième est une exécution : quand elle
se termine, le graphe existe déjà. Les fonctions du vocabulaire écrivent dans une
variable globale posée par le \texttt{with} courant --- d'où l'importance de
l'indentation. Le post-traitement résout ensuite les chemins et les dépendances,
et c'est là que la plupart des erreurs se manifestent, loin de leur cause. Le
namespace du fichier racine est propagé aux inclusions, ce qui rend un
vocabulaire maison possible mais rend aussi son origine invisible. Enfin, la
ligne \texttt{Entry file} de \texttt{jenga info} dit quel workspace a réellement
été chargé.
\end{nkbilan}

\begin{nkquestions}
\begin{enumerate}
\item À quel moment les projets existent-ils en mémoire ?
\item Pourquoi \texttt{utf-8-sig} et non \texttt{utf-8} ?
\item Que fait le \texttt{finally} du chargeur, et qu'éviterait-on sans lui ?
\item D'où vient une fonction employée dans un \texttt{.jenga} sans y être
      importée ?
\item Quelle ligne de \texttt{jenga info} dit quel fichier a été chargé ?
\end{enumerate}
\end{nkquestions}

\begin{nkpratique}
Placez un \texttt{print} au début, au milieu et à la fin de votre
\texttt{.jenga}, ainsi qu'au début d'un fichier inclus.

Lancez \texttt{jenga info} et notez l'\textbf{ordre} d'apparition. Puis déplacez
l'\texttt{include} avant les définitions dont il dépend, et observez.

Retirez tous les \texttt{print} ensuite : ils s'exécutent à \emph{chaque}
commande, y compris celles dont un script lit la sortie.
\end{nkpratique}

\begin{nkdemo}
Prouvez que l'\texttt{include} s'exécute \emph{en ligne}, et non après.

Dans le fichier racine, définissez une variable \texttt{MARQUE = "A"}. Placez un
\texttt{include}. Après lui, changez \texttt{MARQUE = "B"}. Dans le fichier
inclus, affichez \texttt{MARQUE}.

\textbf{Ce que la démonstration établit} : le fichier inclus affiche \texttt{A}.
Il voit donc l'état du namespace \emph{au moment où il est inclus}, pas son état
final.

Déplacez ensuite l'\texttt{include} après l'affectation à \texttt{"B"} et
refaites la mesure. C'est la paire qui prouve, pas le premier essai.
\end{nkdemo}

\begin{nklogique}
Vous lancez \texttt{jenga build} depuis un sous-dossier de votre dépôt et obtenez
« projet inconnu », alors que la même commande fonctionne depuis la racine.

\begin{enumerate}
\item Énumérez deux explications --- l'une sur le fichier chargé, l'autre sur ce
      qu'il contient.
\item Quelle commande, et quelle ligne de sa sortie, les distingue en une
      seconde ?
\item Une troisième explication n'a rien à voir avec Jenga : laquelle, et
      qu'est-ce qui la trahirait ?
\end{enumerate}
\end{nklogique}
